\documentclass{article}

\usepackage{float}
\restylefloat{table}

\usepackage{booktabs}

\title{Team Productivity: POC\\\progname}

\author{\authname}

\date{}

\input{../Comments.text}
\input{../Common.text}

\begin{document}

\maketitle

This document summarizes the contributions of each team member up to the POC
Demo.  The time period of interest is the time between the beginning of the term
and the POC demo.

% \wss{Please do not delete sections or modify the table formats.  Standardization
% helps the instructors and TAs to review these documents. You may add sections to
% the report.}

% \wss{Remove the instructor guidelines/comments in the final version.}

% \wss{Fill in the requested information, like project url, TA name, supervisor name, etc.}

\section{Demo Plans}

\wss{What is your feasibility risk?}

\wss{What will you be demonstrating}

\section{Team Meeting Attendance}

% \wss{For each team member how many team meetings have they attended over the
% time period of interest.  This number should be determined from the meeting
% issues in the team's repo.  The first entry in the table should be the total
% number of team meetings held by the team.}

\begin{table}[H]
\centering
\begin{tabular}{ll}
\toprule
\textbf{Student} & \textbf{Meetings}\\
\midrule
Total & 16\\
Omar Ghoneim & 16\\
Meena Al Omran & 16\\
Marko Popovic & 16\\
Maya Azar & 16\\
Amirah Abdelrahman & 16\\
\bottomrule
\end{tabular}
\end{table}

Team meetings and attendance are tracked through an internal, shared document on Teams. In addition, meetings are scheduled on Outlook calendar, so this can be referred back to when needed.

\section{Supervisor/Stakeholder Meeting Attendance}

% \wss{For each team member how many supervisor/stakeholder team meetings have
% they attended over the time period of interest.  This number should be determined
% from the supervisor meeting issues in the team's repo.  The first entry in the
% table should be the total number of supervisor and team meetings held by the
% team.  If there is no supervisor, there will usually be meetings with
% stakeholders (potential users) that can serve a similar purpose.}

\noindent \textbf{Supervisor's Name: } Dr. Mayu Nishimura

\begin{table}[H]
\centering
\begin{tabular}{ll}
\toprule
\textbf{Student} & \textbf{Meetings}\\
\midrule
Total & 3\\
Omar Ghoneim & 3\\
Meena Al Omran & 3\\
Marko Popovic & 3\\
Maya Azar & 3\\
Amirah Abdelrahman & 3\\
\bottomrule
\end{tabular}
\end{table}

The team keeps track of supervisor meetings and attendance through an internal, shared document on Teams. The meetings themselves are scheduled through Outlook email exchanges and Outlook calendar, although they typically take place on Zoom.

\section{Lecture Attendance}

% \wss{For each team member how many lectures have they attended over the time
% period of interest.  This number should be determined from the lecture issues in
% the team's repo. You can find the number of lectures in the time period of
% interest by looking at the
% \href{https://calendar.google.com/calendar/u/0/embed?src=rnboqiaki1k2la7rpu3bn0um58@group.calendar.google.com&ctz=America/Toronto}
% {Google calendar} for the capstone course. The first entry in the
% table should be the total number of lectures held for the class during the period of interest.}

% \wss{NOTE: There will be approximately 13 lectures between the start of class
% and the POC demos}

\begin{table}[H]
\centering
\begin{tabular}{ll}
\toprule
\textbf{Student} & \textbf{Lectures}\\
\midrule
Total & 11\\
Omar Ghoneim & 5\\
Meena Al Omran & 5\\
Marko Popovic & 5\\
Maya Azar & 5\\
Amirah Abdelrahman & 5\\
\bottomrule
\end{tabular}
\end{table}

All team members were presented for the first 3 lectures.
Following team formation, the team aimed to have 1-2 rotating members present for each lecture who then updated the team about important information, to accomodate
commitments and commutes.
Note that each team member tracked their individual lecture attendance on personal documents.

\section{TA Document Discussion Attendance}

% \wss{For each team member how many of the informal document discussion meetings
% with the TA were attended over the time period of interest. The first row is the
% total number of discussion meetings scheduled by the TA with the team.}

\noindent \textbf{TA's Name: } Rashad Bhuiyan

\begin{table}[H]
\centering
\begin{tabular}{ll}
\toprule
\textbf{Student} & \textbf{Meetings}\\
\midrule
Total & 4\\
Omar Ghoneim & 4\\
Meena Al Omran & 4\\
Marko Popovic & 4\\
Maya Azar & 4\\
Amirah Abdelrahman & 3\\
\bottomrule
\end{tabular}
\end{table}

% \wss{If needed, an explanation for the attendance can be provided here.}
Any absences were due to extranuating circumstances such as illness or medical appointments, and were communicated in advance with and approved by Rashad.

\section{Commits, Lines Added and Lines Deleted, Pull Requests}

\begin{table}[H]
\centering
\begin{tabular}{lllllll}
\toprule
\textbf{Student} & \textbf{Commits} & \textbf{Percent} & \textbf{Lines Added} &
\textbf{Lines Deleted} & \textbf{Pull Requests}\\
\midrule
Total & 180 & 100\% & 28,190 & 3,200 & 28\\
Omar Ghoneim & 44 & 24.44\% & 10,318 & 723 & 16\\
Meena Al Omran & 68 & 37.78\% & 13,247 & 1201 & 8\\
Marko Popovic & 16 & 8.89\% & 904 & 289 & 0\\
Maya Azar & 32 & 17.78\% & 2,188 & 789 & 2\\
Amirah Abdelrahman & 20 & 11.11\% & 1,533 & 198 & 2\\
\bottomrule
\end{tabular}
\end{table}

Omar and Meena retain a disproportionately high number of commits, pull requests and lines added that do not reflect the contributions of other team members. This is partially because Omar pushed the initial Godot project to main, which comes with a large amount of boilerplate code. Meena worked extensively to fix LaTeX-related issues, meaning she had to frequently add and edit lines of code to format and repair issues.

All members contributed to the project in a reasonable way that may not be reflected in this table. For example, Marko contributed to the project by creating a large number of assets to the ongoing implementation, Maya has been working on the project management side extensively, while Amirah has done tons of work on research and connecting with external actors such as professors, teachers and potential users.

\section{Issue Tracker}

\wss{For each team member how many issues have they authored (including open and
closed issues (O+C)) and how many have they been assigned (only counting closed
issues (C only)) over the time period of interest.}

\begin{table}[H]
\centering
\begin{tabular}{lll}
\toprule
\textbf{Student} & \textbf{Authored (O+C)} & \textbf{Assigned (C only)}\\
\midrule
Name 1 & Num & Num \\
Name 2 & Num & Num \\
Name 3 & Num & Num \\
Name 4 & Num & Num \\
Name 5 & Num & Num \\
\bottomrule
\end{tabular}
\end{table}

\wss{If needed, an explanation for the counts can be provided here.}

\section{CICD via GitHub Actions}

% \wss{Say how CICD technology is used and will be used in your project.}

% \wss{Provide links to your CICD yaml files}
Currently, CICD technology is implemented into the project through GitHub Actions for the project documentation, which has a complex \LaTeX{} structure. We
are using the template repository's built-in \href{https://github.com/meena-alomran/game-n-grow/blob/main/.github/workflows/latex-pages.yml}{workflow}. Through this workflow,
GitHub Actions triggers CICD whenever a pull request is made or merged to the main branch. When a pull request is made into the main branch, the workflow tries to compile the \LaTeX{} into pdfs
and registers as a failed run if this is unsuccessful. Then, when the pull request is actually merged, the pdfs are uploaded to a website through GitHub Pages for easy viewing. In addition to this, 
the new pdfs generated from the updated \LaTeX{} code are actually commited into the repository, effectively updating the current pdfs with the changes made to documentation in the pull request.\\
As the project source code is still in beginning stages, CICD has not yet been incorporated for this aspect, but it is in the the team's plan to do once the work is more fleshed out so that building and testing are automated for pull requests into the main branch.

\section{Custom Documents (CDs)}

% \wss{What is the plan (as documented in TeamComposition.csv) for the team's
% CDs?  Should the CDs be modified now that the team knows more about the
% project?}
The original plan for the team's CDs was Stakeholder Interviews for the Fall Term, and a Usability Report for the Winter term.
For Fall CD, it was determined by the team that Interviews were the right choice after all. A crucial part of \progname{} is the unique experience and needs of its users, and interviews
are a good way to get insight into these things so that the game experience is optimal and targets real issues while avoiding unintended consequences. The Stakeholder Interviews were completed, documented, and analyzed (\url{https://github.com/meena-alomran/game-n-grow/blob/main/pdfs/CDs/Interviews/InterviewReport.pdf}), and
the results gave the team a lot of insight for the project.\\
As of now, the plan remains to conduct a Usability Report for the Winter term. This choice buils on the earlier sentiment that the core of \progname{} is the user experience, and will provide a way for the team to analyze how real users might interact with \progname's software and help improve the design.

\section{Team Charter Trigger Items}

\wss{Provide a summary of the quantified triggers identified in the team's
charter.}

\wss{Provide a list of any violations of the triggers.  If the team wishes, the
violations can be summarized on aggregate, instead of naming specific team
members.}

\wss{Provide a plan to address the violations.  This could include revising the
triggers, if they are found to be too weak, strong or ambiguous.}

\section{Additional Productivity Metrics}

% \wss{If your team has additional metrics of productivity, please feel free to
% add them to this report.  If not, please explicitly state that there are no additional
% metrics.}

% \wss{Additional metrics can include things like code reviews done, pull requests
% created, count of joining meetings late, count of number of times contributions
% had to be corrected, number of internal deadlines missed, test cases written, etc.}

% \wss{We are looking for data on these metrics, not just a list of additional
% metrics the team is planning on tracking.  However, if all you have is a plan,
% please share it here.}

There are no additional productivity metrics for this project.

\end{document}